\subsection{Two failure mechanisms outside the window}\label{sec:loadaxis}
Below the critical value, the fluid trajectory of peeling leaves a positive gap at the bottleneck, so that forming a
macroscopic core requires a large fluctuation, whose probability decays rapidly as the load decreases. At the same time
there is another failure mechanism almost independent of the load position: two elements choose exactly the same
$k$-subset. The two edges block each other, each endpoint has degree at least two, and peeling can never remove them,
so they constitute a core consisting of exactly two edges. Its probability is about $\lambda = \binom{d}{2}/\binom{M}{k}$,
which is small in the sparse low-load region, but the probability of a macroscopic failure decays even faster and is
eventually overtaken by it. At a fixed subcritical load, failure is almost always caused by two repeated edges
(Proposition~\ref{prop:lowload} of Appendix~\ref{app:lowload}).

The effect of a repeated edge on the measurement can be computed directly. If the two repeated edges have the same
sign, their counts on the $k$ shared cells add up, the energy exceeds that of two non-overlapping edges by about $2k$,
and the estimate is shifted by about $+2$; if they have opposite signs, the counts cancel and the estimate is shifted by
about $-2$. Conditioned on the presence of a repeated edge, the probability that the two edges have the same sign is
$(1+\gamma)/2$, where $\gamma$ is the sign factor of Section~\ref{sec:sign}. On average the estimate is shifted by
$2\gamma$. This shift is constant on the original scale, while the standard deviation $s$ grows with the system size, so
the standardized measurement at the repeated-edge layer tends again to $N(0,1)$.

Between the critical window and a fixed subcritical load lies the moderate deviation region ($r \to -\infty$ and
$|r| \le C\sqrt{\log M}$), where the two mechanisms coexist. Proposition~\ref{prop:mixture}
(Appendix~\ref{app:mixture}) proves that failures caused by a macroscopic core still obey the selection law, and the law
of the measurement conditioned on failure is a weighted mixture of the macroscopic layer and the repeated-edge layer,
with the two weights proportional to $\Phi(z_r)$ and $\lambda$, respectively. As the load moves down from the critical
window, $\Phi(z_r)$ decays like a normal tail while $\lambda$ stays almost unchanged, so the weight passes smoothly from
the macroscopic layer to the repeated-edge layer.

Replacing the null distribution in the candidate test by this mixture law gives an upper bound covering the entire load
axis. For a candidate $n$ write $z_n = z(n)$, $c_n = \eta_k b^2/n^2$,
$\lambda_n = \binom{n}{2}/\binom{M}{k}$, and $w_n = \Phi(z_n)/(\Phi(z_n) + \lambda_n)$, and set
\begin{equation*}
H_n(x) = w_n\frac{\Phi_2(x,z_n;c_n)}{\Phi(z_n)} + (1-w_n)\Phi(x),\qquad
T_n = n + s(n)\,H_n^{-1}(u).
\end{equation*}
Here $H_n$ is the mixture CDF of the macroscopic layer (weight $w_n$, selection law) and the repeated-edge layer
(weight $1-w_n$, standard normal). Replacing $t(n,b)$ of Section~\ref{sec:upper-construct} by $T_n$ and leaving the
rest of the construction unchanged gives an upper bound, still denoted by $U$.

\begin{proposition}[Load-axis coverage]\label{prop:axis}
Fix $k \ge 3$ and $u \in (0,1)$. If $d/\sqrt{M} \to \infty$, then, whatever the limit of $d/M$ is, whichever side of
the critical point the load lies on, and whether or not the sign composition converges, we have
\begin{equation*}
\sup_x\bigl|\Pr(X \le x \mid F) - H_d(x)\bigr| \to 0,\qquad
\liminf_M \Pr(d \le U \mid F) \ge 1-u.
\end{equation*}
The condition $d/\sqrt{M} \to \infty$ cannot be dropped. Take balanced signs and an even $d \to \infty$ with
$d = o(\sqrt{M})$; then for $u \in (0,1/2)$ we have $\Pr(U < d \mid F) \to 1/2$. For fixed $d = 4$ with two edges of
each sign, this probability tends to $2/3$.
\end{proposition}
The counterexample comes from the $\pm 2$ shift of repeated edges. The order of magnitude of $s$ is $d/\sqrt{M}$, so
when $d = o(\sqrt{M})$ the shift $2/s$ does not vanish on the standardized scale; opposite-sign repeated edges push the
estimate too low, and the true value is then rejected. The proof and the construction of the counterexample are given in
Appendix~\ref{app:axis}.
